% DERIVED from formal_layer.json — read-only, do not edit.
\begin{assumptionv}{ass:iid-sampling}[Independent and identically distributed observations]
For each \(\mathbb P\in\mathcal D_{d,\epsilon}^{\mathrm{obs}}\), the observed-law class, the observed units \(O_1,\ldots,O_n\) are independent and have common law \(\mathbb P\).
\end{assumptionv}

\begin{definitionv}{synth_1}[Covariate alphabet and atoms]
The finite covariate alphabet \([d]\) and observed atom index set \(\mathcal J_d\) are defined, for each nonnegative integer \(d\), by
\[
[d] := \{x\in\mathbb Z : 0\le x<d\},
\qquad
\mathcal J_d := [d]\times\{0,1\}\times\{0,1\}.
\]
\end{definitionv}

\begin{assumptionv}{ass:shrinking-overlap}[Public treatment overlap restriction]
Let \(p_x\) denote the probability of covariate cell \(x\), and let \(\pi_x\) denote its conditional treatment-one probability. For every \(x\in[d]\), the finite covariate alphabet, with \(p_x>0\), the overlap parameter \(\epsilon\) satisfies
\[
\epsilon\le\pi_x\le1-\epsilon.
\]
\end{assumptionv}

\begin{definitionv}{def:observed-class}[Observed laws under overlap]
The observed class \(\mathcal D_{d,\epsilon}^{\mathrm{obs}}\) consists of probability laws \(\mathbb P\) on \([d]\times\{0,1\}^2\) satisfying \cref{obj:ass:shrinking-overlap} on every cell with \(p_x>0\). Its parameters are \(d\in\mathbb N\) and \(\epsilon\in\mathbb R\), and the covariate alphabet is
\[
[d]:=\{x\in\mathbb Z:0\le x<d\}.
\]
For the coordinate variables \(X\), \(A\), and \(Y\), representing the covariate, treatment, and observed outcome, respectively, the observed atom probabilities, cell probabilities, treatment propensities, and arm-specific outcome means are defined by
\[
q_{ay,x}:=\mathbb P(X=x,A=a,Y=y),
\qquad
p_x:=\sum_{a\in\{0,1\}}\sum_{y\in\{0,1\}}q_{ay,x},
\]
\[
\pi_x:=\frac{q_{10,x}+q_{11,x}}{p_x},
\qquad
\mu_{ax}:=\frac{q_{a1,x}}{q_{a0,x}+q_{a1,x}},
\]
with each ratio defined as zero when its denominator is zero. These quantities satisfy, for every \(x\in[d]\) and \(a\in\{0,1\}\),
\[
q_{a1,x}=p_x\pi_x^a(1-\pi_x)^{1-a}\mu_{ax},
\qquad
q_{a0,x}=p_x\pi_x^a(1-\pi_x)^{1-a}(1-\mu_{ax}),
\]
\[
p_x\ge0,
\qquad
\sum_{x\in[d]}p_x=1,
\qquad
\mu_{ax}\in[0,1].
\]
\end{definitionv}

\begin{definitionv}{def:observed-value}[Observed optimal treatment value]
The observed value of a probability law \(\mathbb P\) on \([d]\times\{0,1\}\times\{0,1\}\), for a nonnegative integer \(d\) and the zero-based covariate alphabet \([d]\) of \cref{obj:synth_1}, is
\[
\Psi(\mathbb P)
=\sum_{x\in[d]}p_x\max_{a\in\{0,1\}}\mu_{ax}.
\]
Here the observed atom probabilities, covariate cell probabilities, and arm-specific outcome means are defined by
\[
q_{ay,x}=\mathbb P(X=x,A=a,Y=y),
\qquad
p_x=\sum_{a\in\{0,1\}}\sum_{y\in\{0,1\}}q_{ay,x},
\qquad
\mu_{ax}=\frac{q_{a1,x}}{q_{a0,x}+q_{a1,x}},
\]
with a zero denominator giving \(\mu_{ax}=0\). Each null-cell summand is therefore zero.
\end{definitionv}

\begin{definitionv}{def:minimax-risk}[Minimax risk \(\mathfrak R_{n,d,\epsilon}\)]
The minimax squared-error risk is
\[
\mathfrak R_{n,d,\epsilon}
:=
\inf_{\widehat V}
\sup_{\mathbb P\in\mathcal D_{d,\epsilon}^{\mathrm{obs}}}
E_{\mathbb P}
\left[
\left\{
\widehat V(O_1,\ldots,O_n)-\Psi(\mathbb P)
\right\}^{2}
\right],
\]
where \(O_1,\ldots,O_n\) denotes the observed sample, \(\Psi(\mathbb P)\) denotes the observed value, and the infimum ranges over all measurable real-valued estimators.
\end{definitionv}

\begin{assumptionv}{ass:consistency}[Consistency of observed outcomes]
The observed outcome \(Y\) equals the potential outcome \(Y(A)\) under the realized treatment \(A\) almost surely:
\[
Y=Y(A)\quad\text{almost surely}.
\]
\end{assumptionv}

\begin{assumptionv}{ass:conditional-exchangeability}[Conditional exchangeability of treatment]
For each treatment \(a\in\{0,1\}\), the potential outcome \(Y(a)\) is conditionally independent of the treatment \(A\) given the categorical covariate \(X\):
\[
Y(a)\perp A\mid X.
\]
\end{assumptionv}

\begin{definitionv}{def:causal-class}[Causal-law class \(\mathcal P_{d,\epsilon}\)]
Define
\[
\mathcal P_{d,\epsilon}
:=
\left\{
P:
P\text{ is a law of }(X,A,Y,Y(0),Y(1))
\text{ satisfying the conditions below}
\right\}.
\]
Here \(\operatorname{Obs}(P)\) denotes the observed marginal law of \((X,A,Y)\).
\begin{itemize}
\item \textbf{(Consistency.)} \(P\) satisfies \cref{obj:ass:consistency}.
\item \textbf{(Exchangeability.)} \(P\) satisfies \cref{obj:ass:conditional-exchangeability}.
\item \textbf{(Observed marginal.)}
\[
\operatorname{Obs}(P)\in\mathcal D_{d,\epsilon}^{\mathrm{obs}},
\]
whose overlap restriction is given in \cref{obj:ass:shrinking-overlap}.
\end{itemize}
This definition applies for every \(n\ge1\), \(d\ge2\), and \(0<\epsilon\le1/2\).
\end{definitionv}

\begin{definitionv}{def:armwise-extension}[Armwise extension \(F_\epsilon\)]
The armwise extension is
\[
F_\epsilon(u)
=
\begin{cases}
\displaystyle
\max_{a\in\{0,1\}}\frac{s(u)u_{a1}}{D_a(u)},&u\ne0,\\[6pt]
0,&u=0,
\end{cases}
\qquad u\in\mathbb R_+^4,
\]
where the coordinates of \(u\) are indexed by \(a,y\in\{0,1\}\), and
\[
s_a(u)=u_{a0}+u_{a1},
\qquad
s(u)=s_0(u)+s_1(u),
\qquad
D_a(u)=\max\{s_a(u),\epsilon s(u)\}.
\]
\end{definitionv}

\begin{propositionv}{prop:identification}[Completion and value identification]
Let \(d\) be the number of covariate cells and \(\epsilon\) the overlap floor, with
\begin{itemize}
\item \textbf{(Cell count.)} \(d\in\mathbb N\) and \(d\ge 2\).
\item \textbf{(Positive overlap.)} \(0<\epsilon\).
\item \textbf{(Overlap range.)} \(\epsilon\le 1/2\).
\end{itemize}
The following statements hold.

Every observed law \(\mathbb P\in\mathcal D_{d,\epsilon}^{\mathrm{obs}}\) of \cref{obj:def:observed-class} admits a full-data completion \(P\in\mathcal P_{d,\epsilon}\) of \cref{obj:def:causal-class} such that
\[
\operatorname{Obs}(P)=\mathbb P,
\]
where \(\operatorname{Obs}(P)\) is the marginal law of the categorical covariate \(X\), binary treatment \(A\), and binary observed outcome \(Y\).

Every \(P\in\mathcal P_{d,\epsilon}\) satisfies
\[
V^\star(P)
=
\sum_{x\in[d]} P(X=x)
\max_{a\in\{0,1\}}
\mathbb E_P[Y(a)\mid X=x]
=
\Psi\{\operatorname{Obs}(P)\},
\]
where \([d]\) is the finite covariate alphabet, \(Y(a)\) is the potential outcome under arm \(a\), \(V^\star(P)\) is the causal oracle value defined by the displayed sum, and \(\Psi\) is the observed-value functional of \cref{obj:def:observed-value}. Conditional potential-outcome means in cells of zero probability are assigned zero.

For every \(\mathbb P\in\mathcal D_{d,\epsilon}^{\mathrm{obs}}\) and every \(x\in[d]\), define the four-vector of observed atom probabilities and its cell mass and armwise outcome means by
\[
q_{ay,x}=\mathbb P(X=x,A=a,Y=y),
\qquad
q_x=(q_{00,x},q_{01,x},q_{10,x},q_{11,x}),
\]
\[
p_x=\sum_{a=0}^{1}\sum_{y=0}^{1}q_{ay,x},
\qquad
\mu_{ax}=\frac{q_{a1,x}}{q_{a0,x}+q_{a1,x}},
\]
with a zero denominator giving \(\mu_{ax}=0\). The armwise extension of \cref{obj:def:armwise-extension} then satisfies
\[
F_\epsilon(q_x)=p_x\max_{a\in\{0,1\}}\mu_{ax}.
\]
In particular, both sides equal zero when \(p_x=0\).
\end{propositionv}

\begin{theoremv}{thm:matched-frontier}[Universal matched risk frontier]
There exist universal real constants \(c_\star,C_\star,H_0,\kappa\) and an integer \(D_0\) satisfying
\[
0<c_\star\le C_\star,\qquad H_0>0,\qquad \kappa>0,\qquad D_0\ge2,
\]
such that the following bounds hold for every:
\begin{itemize}
\item \textbf{(Sample size.)} Integer \(n\ge1\).
\item \textbf{(Alphabet size.)} Integer \(d\ge2\).
\item \textbf{(Overlap.)} Real overlap parameter \(0<\epsilon\le1/2\).
\end{itemize}
Write \(L_d=\log(ed)\) for the logarithmic alphabet scale, and let
\(\widehat V_{n,d,\epsilon}^{\mathrm{AF}}\) be the prescribed deterministic armwise estimator with calibration constants \(H_0,\kappa,D_0\). The observed minimax squared-error risk,
\[
\mathfrak R_{n,d,\epsilon}
=
\inf_{\widehat V}
\sup_{\mathbb P\in\mathcal D_{d,\epsilon}^{\mathrm{obs}}}
E_{\mathbb P}\!\left[
\bigl(\widehat V-\Psi(\mathbb P)\bigr)^2
\right],
\]
where the infimum ranges over measurable real-valued estimators based on \(n\) observed units, satisfies
\[
c_\star\min\left\{1,\frac{d}{n\epsilon L_d}\right\}
\le
\mathfrak R_{n,d,\epsilon}
\le
\sup_{\mathbb P\in\mathcal D_{d,\epsilon}^{\mathrm{obs}}}
E_{\mathbb P}\!\left[
\bigl(\widehat V_{n,d,\epsilon}^{\mathrm{AF}}-\Psi(\mathbb P)\bigr)^2
\right]
\le
C_\star\min\left\{1,\frac{d}{n\epsilon L_d}\right\}.
\]
For the causal class, the same estimator and constants satisfy
\[
\begin{aligned}
c_\star\min\left\{1,\frac{d}{n\epsilon L_d}\right\}
&\le
\inf_{\widehat V}
\sup_{P\in\mathcal P_{d,\epsilon}}
E_{\operatorname{Obs}(P)}\!\left[
\bigl(\widehat V-V^\star(P)\bigr)^2
\right]
\\
&\le
\sup_{P\in\mathcal P_{d,\epsilon}}
E_{\operatorname{Obs}(P)}\!\left[
\bigl(\widehat V_{n,d,\epsilon}^{\mathrm{AF}}-V^\star(P)\bigr)^2
\right]
\\
&\le
C_\star\min\left\{1,\frac{d}{n\epsilon L_d}\right\},
\end{aligned}
\]
where the infimum again ranges over measurable real-valued estimators based on \(n\) observed units, \(V^\star(P)\) is the causal oracle value, and the expectation is under the observed marginal law \(\operatorname{Obs}(P)\).
\end{theoremv}

\begin{theoremv}{thm:consistency-threshold}[Consistency threshold and fixed-parameter rates]
Let \(\mathfrak R_{n,d,\epsilon}\) be the minimax squared-error risk in \cref{obj:def:minimax-risk}. Consider sequences \((d_n)_{n\in\mathbb N}\) of alphabet sizes and \((\epsilon_n)_{n\in\mathbb N}\) of overlap floors satisfying:
\begin{itemize}
\item \textbf{(Integer alphabet.)} \(d_n\in\mathbb N\) for every \(n\).
\item \textbf{(Alphabet size.)} \(d_n\ge2\) for every \(n\).
\item \textbf{(Positive overlap.)} \(\epsilon_n>0\) for every \(n\).
\item \textbf{(Overlap range.)} \(\epsilon_n\le1/2\) for every \(n\).
\end{itemize}
There exists a sequence of measurable real-valued estimators \(\widehat V_n\), each based on \(n\) independent observed units, such that
\[
\sup_{\mathbb P\in\mathcal D_{d_n,\epsilon_n}^{\mathrm{obs}}}
\mathbb E_{\mathbb P}
\bigl[(\widehat V_n-\Psi(\mathbb P))^2\bigr]
\longrightarrow0
\]
if and only if
\[
\frac{d_n}{n\epsilon_n\log(e d_n)}
\longrightarrow0
\qquad (n\to\infty).
\]

For each fixed integer \(d\ge2\), there exist real constants \(c,C\), depending on \(d\), with \(0<c\le C\), such that for every integer \(n\ge1\) and every overlap floor \(0<\epsilon\le1/2\),
\[
c\min\left\{1,\frac{1}{n\epsilon}\right\}
\le \mathfrak R_{n,d,\epsilon}
\le C\min\left\{1,\frac{1}{n\epsilon}\right\}.
\]

For each fixed overlap floor \(0<\epsilon\le1/2\), there exist real constants \(c,C\), depending on \(\epsilon\), with \(0<c\le C\), such that for every pair of integers \(n\ge1\) and \(d\ge2\),
\[
c\min\left\{1,\frac{d}{n\log(e d)}\right\}
\le \mathfrak R_{n,d,\epsilon}
\le C\min\left\{1,\frac{d}{n\log(e d)}\right\}.
\]
\end{theoremv}

\begin{algorithmv}{def:armwise-estimator}[Armwise factorial estimator]
The estimator \(\widehat V_{n,d,\epsilon}^{\mathrm{AF}}\) is a deterministic function of the observed sample \(O_i=(X_i,A_i,Y_i)\in[d]\times\{0,1\}\times\{0,1\}\), \(i=1,\ldots,n\), using the zero-based alphabet \([d]=\{0,\ldots,d-1\}\) of \cref{obj:synth_1}. Its parameters satisfy:
\begin{itemize}
\item \textbf{(Tuning.)} The pilot-radius constant \(H_0\) and degree constant \(\kappa\) satisfy \(H_0>0\) and \(\kappa>0\); the integer alphabet cutoff \(D_0\) satisfies \(D_0\ge2\).
\item \textbf{(Sample and alphabet.)} The integers \(n\) and \(d\) satisfy \(n\ge1\) and \(d\ge2\).
\item \textbf{(Overlap parameter.)} The parameter \(\epsilon\) satisfies \(0<\epsilon\le1/2\).
\end{itemize}
Define the logarithmic scale, evaluation intensity, and approximation degree by
\[
L=L_d=\log(ed),\qquad
m=\frac n8,\qquad
K_n=\max\{2,\lfloor\kappa L\rfloor\}.
\]
Write \(\Pi_{[l,u]}(z)=\min\{u,\max\{l,z\}\}\) for projection onto an interval \([l,u]\).

For \(d<D_0\), define the full-sample counts
\[
N_x=\sum_{i=1}^{n}\mathbf 1\{X_i=x\},\qquad
N_{ax}=\sum_{i=1}^{n}\mathbf 1\{X_i=x,A_i=a\},\qquad
S_{ax}=\sum_{i=1}^{n}Y_i\mathbf 1\{X_i=x,A_i=a\},
\qquad x\in[d],\quad a\in\{0,1\},
\]
and set
\[
\widehat V_{n,d,\epsilon}^{\mathrm{AF}}
=
\Pi_{[0,1]}
\left(
\sum_{x\in[d]}\frac{N_x}{n}
\max_{a\in\{0,1\}}
\left\{\frac{S_{ax}}{N_{ax}}\right\}
\right),
\]
where each ratio \(S_{ax}/N_{ax}\) is defined to be zero when \(N_{ax}=0\).

For \(d\ge D_0\) and \(n\epsilon<d/L\), set
\[
\widehat V_{n,d,\epsilon}^{\mathrm{AF}}=\frac12.
\]

For \(d\ge D_0\) and \(n\epsilon\ge d/L\), define the estimator by the following auxiliary construction. Independently of the observations, draw
\[
\mathsf M\sim\operatorname{Pois}(n/4).
\]
On \(\{\mathsf M\le n\}\), draw independent fair marks, independently of the observations and of \(\mathsf M\),
\[
\mathsf B_i\sim\operatorname{Bernoulli}(1/2),
\qquad i=1,\ldots,\mathsf M.
\]
Define the pilot and evaluation atom counts by
\[
N'_{ay,x}
=
\sum_{i=1}^{\mathsf M}
\mathbf 1\{X_i=x,A_i=a,Y_i=y,\mathsf B_i=1\},
\]
\[
N_{ay,x}
=
\sum_{i=1}^{\mathsf M}
\mathbf 1\{X_i=x,A_i=a,Y_i=y,\mathsf B_i=0\},
\qquad a,y\in\{0,1\},\quad x\in[d].
\]
All subsequent cellwise quantities are defined on \(\{\mathsf M\le n\}\). Set
\[
c_{ay,x}=\frac{N'_{ay,x}}m,\qquad
h_{ay,x}
=
H_0\left\{\sqrt{\frac{c_{ay,x}L}{m}}+\frac Lm\right\},
\]
\[
\ell_{ay,x}=\max\{c_{ay,x}-h_{ay,x},0\},\qquad
u_{ay,x}=c_{ay,x}+h_{ay,x},
\]
\[
b_{ay,x}=\frac{\ell_{ay,x}+u_{ay,x}}2,\qquad
r_{ay,x}=\frac{u_{ay,x}-\ell_{ay,x}}2,
\]
\[
b_x=(b_{ay,x})_{a,y\in\{0,1\}},\qquad
r_x=(r_{ay,x})_{a,y\in\{0,1\}},\qquad
Q_x=\prod_{a,y\in\{0,1\}}[\ell_{ay,x},u_{ay,x}].
\]
Use \(F_\epsilon\), \(s\), and \(D_a\) from \cref{obj:def:armwise-extension}. Define the normalized Jackson kernel by
\[
J_{K_n}(t)
=
c_{K_n}^{\mathrm J}
\left\{\frac{\sin(K_nt/2)}{\sin(t/2)}\right\}^{4},
\qquad
\int_{-\pi}^{\pi}J_{K_n}(t)\,dt=1,
\]
where removable values are filled continuously and \(c_{K_n}^{\mathrm J}>0\) is the unique normalizing constant. With \(\odot\) denoting componentwise multiplication and cosine applied componentwise, define the algebraic polynomial \(P_{x,K_n}\) on \(Q_x\) by
\[
P_{x,K_n}(b_x+r_x\odot\cos\theta)
=
\int_{[-\pi,\pi]^4}
F_\epsilon\!\left(b_x+r_x\odot\cos(\theta+\tau)\right)
\prod_{a,y\in\{0,1\}}J_{K_n}(\tau_{ay})\,d\tau,
\qquad \theta\in\mathbb R^4.
\]
Its degree in each coordinate is at most
\[
D=2(K_n-1).
\]
Define the coefficients \(\beta_{x,\boldsymbol\alpha}\) by the unique expansion
\[
P_{x,K_n}(v)-F_\epsilon(b_x)
=
\sum_{\boldsymbol\alpha\in\{0,\ldots,D\}^4}
\beta_{x,\boldsymbol\alpha}
\prod_{a,y\in\{0,1\}}
\left(\frac{v_{ay}-b_{ay,x}}{r_{ay,x}}\right)^{\alpha_{ay}},
\qquad v\in Q_x.
\]
For nonnegative integers \(N\) and \(t\), define the falling factorial by
\[
(N)_t=\prod_{j=0}^{t-1}(N-j),\qquad (N)_0=1.
\]
For an integer \(k\ge0\) and a real center \(\zeta\), define
\[
U_k(N;\zeta)
=
\sum_{t=0}^{k}\binom kt(-\zeta)^{k-t}\frac{(N)_t}{m^t},
\]
and set
\[
Z_x
=
F_\epsilon(b_x)
+
\sum_{\boldsymbol\alpha\in\{0,\ldots,D\}^4}
\beta_{x,\boldsymbol\alpha}
\prod_{a,y\in\{0,1\}}
\frac{U_{\alpha_{ay}}(N_{ay,x};b_{ay,x})}
{r_{ay,x}^{\alpha_{ay}}}.
\]
Define the armwise clipping scales and clipped cell contributions by
\[
R_{a,x}=r_{a0,x}+r_{a1,x},\qquad
w_{a,x}=1+\frac{s(b_x)}{D_a(b_x)},\qquad
S_x^{\mathrm{aw}}=2\sum_{a=0}^{1}w_{a,x}R_{a,x},
\]
\[
T_x
=
\Pi_{[F_\epsilon(b_x)-d^{1/4}S_x^{\mathrm{aw}},\,
       F_\epsilon(b_x)+d^{1/4}S_x^{\mathrm{aw}}]}(Z_x).
\]
Finally, set
\[
\widetilde V
=
\begin{cases}
\displaystyle\Pi_{[0,1]}\left(\sum_{x\in[d]}T_x\right),
&\mathsf M\le n,\\[4pt]
0,&\mathsf M>n,
\end{cases}
\qquad
\widehat V_{n,d,\epsilon}^{\mathrm{AF}}
=
E[\widetilde V\mid O_1,\ldots,O_n].
\]
The conditional expectation integrates the auxiliary count and marks. Equivalently, on this branch,
\[
\widehat V_{n,d,\epsilon}^{\mathrm{AF}}
=
\sum_{k=0}^{n}
e^{-n/4}\frac{(n/4)^k}{k!}
\sum_{(\mathsf B_1,\ldots,\mathsf B_k)\in\{0,1\}^k}
2^{-k}\,
\Pi_{[0,1]}\left(\sum_{x\in[d]}T_x\right),
\]
where each summand uses the counts and cell contributions constructed from the first \(k\) observations and the displayed mark vector. This defines a deterministic measurable real-valued function of the observed sample.
\end{algorithmv}

\begin{theoremv}{thm:observable-upper}[Uniform observable risk bound]
There exist universal constants \(H_0>0\), \(\kappa>0\), \(0<C_\star<\infty\), and an integer \(D_0\ge2\) such that, for every choice of:
\begin{itemize}
\item \textbf{(Sample size.)} An integer \(n\ge1\).
\item \textbf{(Alphabet size.)} An integer \(d\ge2\).
\item \textbf{(Overlap.)} A public overlap parameter \(0<\epsilon\le1/2\).
\end{itemize}
the estimator \(\widehat V_{n,d,\epsilon}^{\mathrm{AF}}\) specified in \cref{obj:def:armwise-estimator}, with these tuning constants, is a measurable, deterministic function of the observed sample. For every family of sampling laws satisfying \cref{obj:ass:iid-sampling} for each \(\mathbb P\in\mathcal D_{d,\epsilon}^{\mathrm{obs}}\), the class of admissible observed laws with public overlap floor \(\epsilon\), its squared-error risk for the observed value \(\Psi(\mathbb P)\) defined in \cref{obj:def:observed-value} satisfies
\[
\sup_{\mathbb P\in\mathcal D_{d,\epsilon}^{\mathrm{obs}}}
E_{\mathbb P}\!\left[
\left\{\widehat V_{n,d,\epsilon}^{\mathrm{AF}}-\Psi(\mathbb P)\right\}^{2}
\right]
\le C_\star\min\left\{1,\frac{d}{n\epsilon L_d}\right\},
\qquad L_d=\log(ed).
\]
\end{theoremv}

\begin{remarkv}{oeq:sharp-path-constant}[Open questions on path constants]
Consider admissible nonsaturated sequences satisfying
\begin{itemize}
\item \textbf{(Growing alphabet.)} \(d\to\infty\).
\item \textbf{(Vanishing overlap.)} \(\epsilon\to0\).
\item \textbf{(Nonsaturation.)} \(d/(n\epsilon L_d)\to0\), where \(L_d\) is the logarithmic alphabet scale.
\end{itemize}
For \(\mathfrak R_{n,d,\epsilon}\), the minimax squared-error risk, consider the normalized sequence
\[
\frac{n\epsilon L_d}{d}\,\mathfrak R_{n,d,\epsilon}.
\]
It remains open whether there exists a single finite positive constant to which the entire normalized-risk sequence converges along every admissible path. Alternatively, do there exist two admissible overlap paths along which the entire normalized-risk sequences converge to distinct finite positive constants? Failure of convergence along an admissible path is a further possibility.

Future work could investigate a cone-constrained approximation--moment saddle for the boundary-propensity likelihood and an observable no-loss transfer.
\end{remarkv}

\begin{lemmav}{lem:armwise-modulus}[Armwise perturbation and mass bounds]
Let \(\epsilon\) be the overlap parameter, with \(0<\epsilon\leq 1/2\). For every pair of four-vectors \(u,v\) satisfying
\begin{itemize}
\item \textbf{(Nonnegativity.)} \(u,v\in\mathbb R_+^4\).
\item \textbf{(Nonzero comparison.)} \(v\neq 0\).
\end{itemize}
define the armwise perturbations by
\[
\Delta_a=\sum_{y=0}^1|u_{ay}-v_{ay}|,\qquad a\in\{0,1\}.
\]
With \(F_\epsilon\), \(s\), and \(D_a\) as in \cref{obj:def:armwise-extension},
\[
|F_\epsilon(u)-F_\epsilon(v)|
\leq
2\sum_{a=0}^1\Delta_a
+
2\sum_{a=0}^1\frac{s(v)}{D_a(v)}\Delta_a.
\]

For every observed law \(\mathbb P\) on
\(\{0,\ldots,d-1\}\times\{0,1\}\times\{0,1\}\) and every cell \(x\in\{0,\ldots,d-1\}\), let \(v=q_x\) be its vector of observed atom probabilities and \(p_x\) its cell probability:
\[
v_{ay}=\mathbb P\bigl(\{(x,a,y)\}\bigr),
\qquad
p_x=\sum_{a=0}^1\sum_{y=0}^1v_{ay}.
\]
If
\begin{itemize}
\item \textbf{(Dimension.)} \(d\) is an integer with \(d\geq2\).
\item \textbf{(Observed class.)} \(\mathbb P\in\mathcal D_{d,\epsilon}^{\mathrm{obs}}\), as in \cref{obj:def:observed-class}.
\item \textbf{(Positive cell mass.)} \(p_x>0\).
\end{itemize}
then, with \(s_a\) as in \cref{obj:def:armwise-extension},
\[
\sum_{a=0}^1\frac{p_x}{s_a(v)}
\sum_{y=0}^1\sqrt{v_{ay}}
\leq
2\sqrt{2}\sqrt{\frac{p_x}{\epsilon}}.
\]
\end{lemmav}

\begin{definitionv}{def:fixed-sample-factorial}[Centered factorial statistic \(U_{\boldsymbol h}\)]
The centered multinomial factorial statistic is
\[
U_{\boldsymbol h}(N^{\mathrm e};b)
=
\sum_{\boldsymbol t\le\boldsymbol h}
\left\{
\prod_{j\in\mathcal J_d}
\binom{h_j}{t_j}(-b_j)^{h_j-t_j}
\right\}
\frac{\prod_{j\in\mathcal J_d}(N_j^{\mathrm e})_{t_j}}
{(n_{\mathrm e})_{|\boldsymbol t|}},
\qquad
|\boldsymbol h|\le n_{\mathrm e}.
\]
Here \(\mathcal J_d\) is the observed atom index set, \(b=(b_j)_{j\in\mathcal J_d}\) is the centering vector, and \(\boldsymbol h,\boldsymbol t\) are nonnegative integer multi-indices, with componentwise ordering and total orders
\[
|\boldsymbol h|=\sum_{j\in\mathcal J_d}h_j,
\qquad
|\boldsymbol t|=\sum_{j\in\mathcal J_d}t_j.
\]
The deterministic sample split and its atom counts are
\[
n_{\mathrm p}=\lfloor n/2\rfloor,
\qquad
n_{\mathrm e}=n-n_{\mathrm p},
\qquad
I_{\mathrm p}=\{1,\ldots,n_{\mathrm p}\},
\qquad
I_{\mathrm e}=\{n_{\mathrm p}+1,\ldots,n\},
\]
\[
N_j^{\mathrm p}
=\sum_{i\in I_{\mathrm p}}\mathbf 1\{O_i=j\},
\qquad
N_j^{\mathrm e}
=\sum_{i\in I_{\mathrm e}}\mathbf 1\{O_i=j\},
\qquad
N^{\mathrm e}=(N_j^{\mathrm e})_{j\in\mathcal J_d}.
\]
The falling factorials are
\[
(m)_k=m(m-1)\cdots(m-k+1),
\qquad
(m)_0=1.
\]
\end{definitionv}

\begin{definitionv}{def:weighted-error}[Weighted approximation error \(\mathcal E_{K,M,\epsilon}^{w}\)]
The weighted best uniform polynomial approximation error is
\[
\mathcal E_{K,M,\epsilon}^{w}
:=
\inf_{P\in\mathbb R_K[z]}
\sup_{0\le z\le M}
\frac{|\phi_\epsilon(z)-P(z)|}{1+z},
\]
where \(\phi_\epsilon\) is the weighted scalar positive-part functional and \(\mathbb R_K[z]\) is the set of real polynomials in \(z\) of degree at most \(K\).
\end{definitionv}

\begin{definitionv}{def:constrained-prior-separation}[Constrained prior separation \(\Delta_{K,M,\epsilon}\)]
The normalized constrained prior separation is
\[
\Delta_{K,M,\epsilon}
=
\sup_{(\nu_0,\nu_1)\in\mathfrak N_{K,M}}
\left|
\int_0^M\phi_\epsilon(z)\,d(\nu_1-\nu_0)(z)
\right|,
\]
where \(\phi_\epsilon\) is the weighted scalar positive-part functional and the parameters satisfy
\begin{itemize}
\item \textbf{(Degree.)} \(K\) is an integer with \(K\geq2\).
\item \textbf{(Support endpoint.)} \(M\) is real with \(2\leq M\leq K^2\).
\item \textbf{(Overlap.)} \(0\leq\epsilon\leq1/2\).
\end{itemize}
The set \(\mathfrak N_{K,M}\) consists of pairs \((\nu_0,\nu_1)\) of probability measures on \([0,M]\) satisfying
\[
\int_0^M z\,d\nu_0(z)
=
\int_0^M z\,d\nu_1(z)
=1
\]
and
\[
\int_0^M z^j\,d\nu_0(z)
=
\int_0^M z^j\,d\nu_1(z),
\qquad j=0,1,\ldots,K.
\]
\end{definitionv}

\begin{definitionv}{def:weighted-handle}[Common-atom completion map \(\mathscr H^w\)]
The common-atom completion map is
\[
\mathscr H^w(\sigma_+,\sigma_-)
=
\left(
\sigma_-+(1-t)_+\delta_{z_0},
\ \sigma_++(1-t)_+\delta_{z_0}
\right),
\]
with
\[
t=\sigma_+(\mathbb R),
\qquad
u=\int_{\mathbb R}z\,d\sigma_+(z),
\qquad
z_0=\frac{1-u}{1-t}.
\]
Its inputs satisfy the following conditions:
\begin{itemize}
\item \textbf{(Measures.)} \(\sigma_+\) and \(\sigma_-\) are finite nonnegative measures on \(\mathbb R\).
\item \textbf{(First moments.)} Each measure has a finite first moment.
\item \textbf{(Mass.)} \(t<1\).
\end{itemize}
Here \((r)_+=\max\{r,0\}\), and \(\delta_{z_0}\) is the unit point mass at \(z_0\).
\end{definitionv}

\begin{lemmav}{lem:jackson-packet-localization}[Jackson packet localization bounds]
There exist real constants \(0<c\leq C\) such that the following holds for every integer \(N\geq2\). Define the normalized periodic Jackson function by
\[
\mathcal J_N(u)
=c_N\left\{\frac{\sin(Nu/2)}{\sin(u/2)}\right\}^{4},
\qquad
c_N=
\left(
\int_{-\pi}^{\pi}
\left\{\frac{\sin(Nu/2)}{\sin(u/2)}\right\}^{4}
\frac{du}{2\pi}
\right)^{-1},
\]
with removable values filled continuously. Set
\[
L=4N,\qquad
k_N(u)=\mathcal J_N(u)\cos(Lu),
\qquad
\widehat h(j)=\int_{-\pi}^{\pi}h(u)e^{-iju}\,\frac{du}{2\pi}.
\]
For \(j\in\mathbb Z\), define the zero-extended coefficient sequences and their forward differences by
\[
a_j^\pm=
\begin{cases}
\widehat{\mathcal J_N}(j)/(L\pm j)^2,& |j|\leq2N-2,\\
0,& |j|>2N-2,
\end{cases}
\qquad
(\delta a)_j=a_{j+1}-a_j,
\qquad
(\delta^2a)_j=a_{j+2}-2a_{j+1}+a_j.
\]
Define also
\[
G_N(u)
=-\frac12\sum_{j=-(2N-2)}^{2N-2}
\left\{
a_j^+\cos((L+j)u)+a_j^-\cos((L-j)u)
\right\},
\qquad
\|u\|_{\mathbb T}
=\inf_{j\in\mathbb Z}|u-2\pi j|.
\]
Then
\[
\int_{-\pi}^{\pi}\mathcal J_N(u)\,\frac{du}{2\pi}=1.
\]
For each sign, the second forward differences are absolutely summable and satisfy
\[
\sum_{j\in\mathbb Z}|\delta^2a_j^\pm|
\leq\frac{C}{N^3}.
\]
The differences include the endpoints of the displayed zero extension. For every \(r\in\mathbb Z\) satisfying
\[
|r|<2N+2\quad\text{or}\quad |r|>6N-2,
\]
the trigonometric integrals vanish:
\[
\int_{-\pi}^{\pi}k_N(u)\cos(ru)\,du=0,
\qquad
\int_{-\pi}^{\pi}k_N(u)\sin(ru)\,du=0.
\]
The function \(G_N\) satisfies
\[
G_N''(u)=k_N(u)\quad\text{for every }u\in\mathbb R,
\qquad
\int_{-\pi}^{\pi}G_N(u)\,du=0,
\]
and the bounds
\[
\frac{c}{N}\leq |G_N(0)|\leq\frac{C}{N},
\qquad
|G_N(u)|\leq
\frac{C}{N\{1+N\|u\|_{\mathbb T}\}^{2}}
\quad\text{for every }u\in\mathbb R,
\qquad
\int_{-\pi}^{\pi}|G_N(u)|\,\frac{du}{2\pi}
\leq\frac{C}{N^2}.
\]
\end{lemmav}

\begin{theoremv}{thm:weighted-separation-and-frontier}[Weighted separation and minimax frontier]
There exist universal real constants \(0<c<C<\infty\) and an integer \(A_0\geq1\) such that the following approximation and separation guarantees hold for every choice of parameters satisfying:
\begin{itemize}
\item \textbf{(Degree.)} \(K\) is an integer with \(K\geq2\).
\item \textbf{(Support.)} \(M\) is real and \(2\leq M\leq K^2\).
\item \textbf{(Overlap.)} \(0\leq\epsilon\leq1/2\).
\end{itemize}
For the weighted error and constrained-prior separation in
\cref{obj:def:weighted-error,obj:def:constrained-prior-separation}, with scalar functional
\[
\phi_\epsilon(z)=
\frac{((1-\epsilon)+\epsilon z)\max\{z-1,0\}}{1+z},
\]
one has
\[
c\frac{\sqrt M}{K}
\leq \mathcal E_{K,M,\epsilon}^{w}
\leq \Delta_{K,M,\epsilon}
\leq 4\mathcal E_{K,M,\epsilon}^{w}
\leq C\frac{\sqrt M}{K}.
\]
Moreover, let \(\sigma_+\) and \(\sigma_-\) be the positive and negative Jordan parts of the normalized Jackson packet with calibration \(A_0\), pushed to \([0,M]\). These are finite measures with integrable first moments. Their positive mass and common completion location satisfy
\[
t=\sigma_+(\mathbb R)\leq\frac12,
\qquad
u=\int z\,d\sigma_+(z),
\qquad
z_0=\frac{1-u}{1-t},
\qquad
\frac12\leq z_0\leq2.
\]
There exists \((\nu_0,\nu_1)\in\mathfrak N_{K,M}\), the constrained prior class of \cref{obj:def:constrained-prior-separation}, such that the common-atom completion of \cref{obj:def:weighted-handle} satisfies
\[
\mathscr H^w(\sigma_+,\sigma_-)
=(\nu_0,\nu_1)
=\bigl(\sigma_-+(1-t)\delta_{z_0},
       \sigma_++(1-t)\delta_{z_0}\bigr),
\]
and
\[
c\frac{\sqrt M}{K}
\leq
\left|
\int\phi_\epsilon(z)\,d\nu_1(z)
-
\int\phi_\epsilon(z)\,d\nu_0(z)
\right|.
\]

There also exist universal real constants
\(0<c_\star\leq C_\star<\infty\), \(H_0>0\), and \(\kappa>0\), and a universal integer \(D_0\geq2\), such that the armwise estimator
\(\widehat V_{n,d,\epsilon}^{\mathrm{AF}}\) with these calibrations satisfies the following risk guarantees for every choice of parameters satisfying:
\begin{itemize}
\item \textbf{(Sample size.)} \(n\) is an integer with \(n\geq1\).
\item \textbf{(Alphabet.)} \(d\) is an integer with \(d\geq2\).
\item \textbf{(Overlap.)} \(0<\epsilon\leq1/2\).
\end{itemize}
The observed risks obey
\[
c_\star\min\left\{1,\frac{d}{n\epsilon\log(ed)}\right\}
\leq \mathfrak R_{n,d,\epsilon}
\leq
\sup_{\mathbb P\in\mathcal D_{d,\epsilon}^{\mathrm{obs}}}
E_{\mathbb P}\!\left[
\left\{\widehat V_{n,d,\epsilon}^{\mathrm{AF}}
-\Psi(\mathbb P)\right\}^2
\right]
\leq
C_\star\min\left\{1,\frac{d}{n\epsilon\log(ed)}\right\}.
\]
The causal minimax risk obeys
\[
c_\star\min\left\{1,\frac{d}{n\epsilon\log(ed)}\right\}
\leq
\inf_{\widehat V}
\sup_{P\in\mathcal P_{d,\epsilon}}
E_{\operatorname{Obs}(P)}\!\left[
\left\{\widehat V-V^\star(P)\right\}^2
\right]
\leq
C_\star\min\left\{1,\frac{d}{n\epsilon\log(ed)}\right\},
\]
where the infimum ranges over measurable real-valued estimators based on \(n\) observed units.

Finally, consider any sequences satisfying:
\begin{itemize}
\item \textbf{(Sample sizes.)} The integer sample sizes \(n_k\) tend to infinity.
\item \textbf{(Alphabets.)} Each \(d_k\) is an integer with \(d_k\geq2\).
\item \textbf{(Overlap.)} \(0<\epsilon_k\leq1/2\) for every \(k\).
\end{itemize}
There exists a sequence of measurable real-valued estimators
\(\widehat V_k\), each based on \(n_k\) observed units with alphabet size \(d_k\), such that
\[
\sup_{\mathbb P\in\mathcal D_{d_k,\epsilon_k}^{\mathrm{obs}}}
E_{\mathbb P}\!\left[
\left\{\widehat V_k-\Psi(\mathbb P)\right\}^2
\right]\longrightarrow0
\]
if and only if
\[
\frac{d_k}{n_k\epsilon_k\log(ed_k)}\longrightarrow0.
\]
\end{theoremv}

\begin{definitionv}{def:sparse-family}[Sparse observed law]
The sparse observed law \(\mathbb P_{\mathbf z}^{\mathrm{sp}}\) is defined for a covariate alphabet size \(d\in\mathbb N\), an overlap parameter \(\epsilon\in\mathbb R\), an intensity support endpoint \(M\in\mathbb R\), and a vector of cell intensities \(\mathbf z\), with
\begin{itemize}
\item \textbf{(Alphabet.)} \(d>0\).
\item \textbf{(Overlap.)} \(0<\epsilon\le 1/2\).
\item \textbf{(Intensities.)} \(\mathbf z\in[0,M]^d\).
\end{itemize}
For each covariate cell \(x\in[d]\), define the unnormalized atom masses by
\[
\widetilde q_{11,x}=\frac{\epsilon(1-\epsilon)z_x}{d},
\qquad
\widetilde q_{10,x}=\frac{\epsilon(1-\epsilon)}{d},
\qquad
\widetilde q_{00,x}=\widetilde q_{01,x}
=\frac{(1-\epsilon)^2+\epsilon^2z_x}{2d}.
\]
Their total mass is
\[
S(\mathbf z)
=\sum_{x\in[d]}\sum_{a=0}^1\sum_{y=0}^1\widetilde q_{ay,x}.
\]
The normalized probability law of the categorical covariate \(X\), binary treatment \(A\), and binary observed outcome \(Y\) is
\[
\mathbb P_{\mathbf z}^{\mathrm{sp}}(X=x,A=a,Y=y)
=\frac{\widetilde q_{ay,x}}{S(\mathbf z)},
\qquad x\in[d],\quad a,y\in\{0,1\}.
\]
\end{definitionv}

\begin{lemmav}{lem:sparse-likelihood}[Sparse law and likelihood identities]
Let the parameters satisfy:
\begin{itemize}
\item \textbf{(Sample and alphabet sizes.)} \(n,d\in\mathbb N\), \(n\ge1\), and \(d\ge2\), with the zero-based cell alphabet \([d]\) of \cref{obj:synth_1}.
\item \textbf{(Overlap.)} \(0<\epsilon\le1/2\).
\item \textbf{(Intensity support.)} \(M\ge2\) and \(\mathbf z\in[0,M]^d\).
\item \textbf{(Poisson inflation.)} \(B>2\).
\end{itemize}
The normalized sparse law \(\mathbb P_{\mathbf z}^{\mathrm{sp}}\) of \cref{obj:def:sparse-family} belongs to \(\mathcal D_{d,\epsilon}^{\mathrm{obs}}\) of \cref{obj:def:observed-class}. For every cell \(x\in[d]\), its cell probability, treatment propensity, and conditional outcome means satisfy
\[
p_x=\frac{1-\epsilon+\epsilon z_x}{dS(\mathbf z)},
\qquad
\pi_x=\frac{\epsilon(1-\epsilon)(1+z_x)}
{1-\epsilon+\epsilon z_x},
\qquad
\mu_{0x}=\frac12,
\qquad
\mu_{1x}=\frac{z_x}{1+z_x}.
\]
Define the scalar functional
\[
\phi_\epsilon(z)
=\frac{(1-\epsilon+\epsilon z)\max\{z-1,0\}}{1+z}.
\]
The observed value of \cref{obj:def:observed-value} is
\[
\Psi(\mathbb P_{\mathbf z}^{\mathrm{sp}})
=\frac12+\frac{1}{2dS(\mathbf z)}
\sum_{x\in[d]}\phi_\epsilon(z_x).
\]

For the one-cell Poisson experiment, write
\(\widetilde q_{ay,x}(z)\) for the unnormalized atom mass in
\cref{obj:def:sparse-family} with cell intensity \(z\), and set the reference intensity \(z_\star=M/2\). Under intensity \(z\), the four counts \(N_{ay,x}\), \(a,y\in\{0,1\}\), are independent Poisson variables with means \(Bn\widetilde q_{ay,x}(z)\). Define their likelihood ratio relative to intensity \(z_\star\) by
\[
L_z(N)
=
\prod_{a,y\in\{0,1\}}
\exp\!\left\{
Bn\bigl(\widetilde q_{ay,x}(z_\star)
-\widetilde q_{ay,x}(z)\bigr)
\right\}
\left(
\frac{\widetilde q_{ay,x}(z)}
{\widetilde q_{ay,x}(z_\star)}
\right)^{N_{ay,x}}.
\]
Writing \(E_{z_\star}\) for expectation under these independent Poisson counts at reference intensity \(z_\star\), for every \(z,w\in[0,M]\),
\[
E_{z_\star}[L_zL_w]
=\exp\!\left\{\alpha(z-z_\star)(w-z_\star)\right\},
\]
where the Gram coefficient is
\[
\alpha
=\frac{Bn}{d}
\left\{
\frac{\epsilon(1-\epsilon)}{z_\star}
+\frac{\epsilon^4}{(1-\epsilon)^2+\epsilon^2z_\star}
\right\}
\le \frac{2Bn\epsilon}{dM}.
\]
\end{lemmav}

\begin{theoremv}{thm:all-estimator-lower}[Universal minimax risk lower bound]
There exists a universal constant \(c_\star>0\) such that, for every sample size \(n\in\mathbb N\), alphabet size \(d\in\mathbb N\), and overlap parameter \(\epsilon\in\mathbb R\) satisfying
\begin{itemize}
\item \textbf{(Sample size.)} \(n\ge1\).
\item \textbf{(Alphabet size.)} \(d\ge2\).
\item \textbf{(Overlap.)} \(0<\epsilon\le1/2\).
\end{itemize}
the minimax squared-error risk \(\mathfrak R_{n,d,\epsilon}\) satisfies
\[
\mathfrak R_{n,d,\epsilon}
\ge c_\star\min\left\{1,\frac{d}{n\epsilon L_d}\right\},
\qquad L_d=\log(ed).
\]
\end{theoremv}
